% Quelle: 4. Übung Lösungen.pdf
% Art: Übung mit Lösungen
% Themen: Inklusion-Exklusion, fixpunktfreie Permutationen, Teileranzahlfunktion, Teilersummenfunktion, vollkommene Zahlen, Mersenne-Primzahlen, kleiner Satz von Fermat, Ordnung, primitive Elemente, Inverse in Zm*, RSA, erweiterter Euklidischer Algorithmus, Square-and-Multiply
\section{Übung 4 (mit Lösungen)}
\subsection{Inklusion-Exklusion und fixpunktfreie Permutationen}

\begin{aufgabe}[Aufgabe 1]
Wie viele ganze Zahlen zwischen 1 und 1000 (einschließlich) sind weder durch 2, noch durch 3 noch durch 5 teilbar?
\end{aufgabe}

\begin{loesung}
% KORRIGIERT: "Omega = {1 <= k <= 1000}" ohne Grundmenge, "A_i = { Omega | i ist Teiler von 1000 }"
$\Omega = \{k \in \Z \mid 1 \le k \le 1000\}$. Sei $A_i = \{\, k \in \Omega \mid i \text{ ist Teiler von } k \,\}$.
Die gesuchte Anzahl ist dann
% KORRIGIERT: Klammern um A_2 u A_3 u A_5 fehlten (hier und unten)
\begin{align*}
|\,\Omega \setminus (A_2 \cup A_3 \cup A_5)\,| &= |\Omega| - \alpha_1 + \alpha_2 - \alpha_3
\end{align*}
mit
\begin{align*}
\alpha_1 &= |A_2| + |A_3| + |A_5|\\
\alpha_2 &= |A_2 \cap A_3| + |A_2 \cap A_5| + |A_3 \cap A_5|\\
\alpha_3 &= |A_2 \cap A_3 \cap A_5|
\end{align*}
\begin{align*}
&|A_2| = \frac{1000}{2}, \qquad |A_3| = \left\lfloor \frac{1000}{3} \right\rfloor, \qquad |A_5| = \frac{1000}{5},\\
\Rightarrow\quad &\alpha_1 = 500 + 333 + 200 = 1033\\
&|A_2 \cap A_3| = \left\lfloor \frac{1000}{6} \right\rfloor, \quad |A_2 \cap A_5| = \frac{1000}{10}, \quad |A_3 \cap A_5| = \left\lfloor \frac{1000}{15} \right\rfloor\\
\Rightarrow\quad &\alpha_2 = 166 + 100 + 66 = 332\\
&|A_2 \cap A_3 \cap A_5| = \left\lfloor \frac{1000}{30} \right\rfloor,\\
\Rightarrow\quad &\alpha_3 = 33\\
\Rightarrow\quad &|\,\Omega \setminus (A_2 \cup A_3 \cup A_5)\,| = |\Omega| - \alpha_1 + \alpha_2 - \alpha_3\\
&= 1000 - 1033 + 332 - 33 = 266
\end{align*}

Schneller findet man die Lösung, wenn man die Produktformel benutzt, die hier aber nur näherungsweise gilt:
% KORRIGIERT: "=" statt "ungefähr gleich"; "266,66" statt 266,6 periodisch
\begin{align*}
|\,\Omega \setminus (A_2 \cup A_3 \cup A_5)\,| &\approx 1000 \cdot \left(1 - \frac{1}{2}\right)\left(1 - \frac{1}{3}\right)\left(1 - \frac{1}{5}\right)\\
&= 100 \cdot \frac{8}{3}\\
&= 266{,}\overline{6}
\end{align*}
% Hinweis: "100 * 8/3" so im Original (entspricht 1000 * 4/15).
\end{loesung}

\begin{aufgabe}[Aufgabe 2]
Wie viele fixpunktfreie Permutationen aus $S_4$ gibt es? Welche sind das?
\end{aufgabe}

\begin{loesung}
\begin{align*}
d_4 &= 4!\left(1 - \frac{1}{1!} + \frac{1}{2!} - \frac{1}{3!} + \frac{1}{4!}\right) = 4!\left(\frac{1}{2} - \frac{1}{6} + \frac{1}{24}\right) = 4! \cdot \frac{9}{24} = 9\\
% KORRIGIERT: fehlender Malpunkt bei 4!; "P = 1/e" ohne Grenzwert
P(4) &= \frac{9}{24} = \frac{3}{8} = 0{,}375 \qquad\qquad \lim_{n \to \infty} P(n) = 1/e = 0{,}3678\ldots
\end{align*}

2 3 4 1, \quad 2 4 1 3, \quad 2 1 4 3, \quad 3 1 4 2, \quad 3 4 2 1, \quad 3 4 1 2,\\
4 1 2 3, \quad 4 3 1 2, \quad 4 3 2 1
\end{loesung}

\subsection{Teileranzahl- und Teilersummenfunktion, vollkommene Zahlen}

% KORRIGIERT: Summationsindex "d" kollidierte mit der Funktion d(n); Voraussetzungen an p_i, e_i fehlten
\[
d(n) = \sum_{t \mid n} 1, \qquad \sigma(n) = \sum_{t \mid n} t \qquad\qquad n = p_1^{e_1} p_2^{e_2} \cdots p_r^{e_r}
\]
mit paarweise verschiedenen Primzahlen $p_1, \ldots, p_r$ und $e_1, \ldots, e_r \geq 1$.

\begin{aufgabe}[Aufgabe 3]
Leiten Sie Formeln für $d(n)$ und $\sigma(n)$ her.\\
Hinweis:
\begin{itemize}
  \item[-] Starten Sie mit dem Fall $n = p$ Primzahl.
  \item[-] Für $\ggT(m, n) = 1$ gilt $d(m\,n) = d(m)\,d(n)$ bzw.\ $\sigma(m\,n) = \sigma(m)\,\sigma(n)$.
\end{itemize}
\end{aufgabe}

\begin{loesung}
Die Teiler von $n = p$ sind $1, p$, also $d(p) = 1 + 1$.\\
Die Teiler von $n = p^e$ sind $1, p, p^2, \ldots, p^e$, also $d(p^e) = 1 + e$
% KORRIGIERT: Begründung (Multiplikativität, ggT = 1, p - 1 != 0) ergänzt; \ldots/\cdots vereinheitlicht
\[
\Rightarrow\qquad d(n) = (1 + e_1) \cdots (1 + e_r) \qquad (\text{da } \ggT(p_i^{e_i}, p_j^{e_j}) = 1 \text{ für } i \neq j)
\]
\[
\sigma(p) = 1 + p, \qquad \sigma(p^e) = 1 + p + p^2 + \cdots + p^e = \frac{p^{1+e} - 1}{p - 1} \quad (p \geq 2, \text{ also } p - 1 \neq 0),
\]
\[
\Rightarrow\qquad \sigma(n) = \frac{p_1^{1+e_1} - 1}{p_1 - 1} \cdots \frac{p_r^{1+e_r} - 1}{p_r - 1}
\]
\end{loesung}

Eine Zahl $n$ mit $\sigma(n) = 2n$ heißt vollkommen.

\begin{aufgabe}[Aufgabe 4]
Zeigen Sie dass 6 und 28 vollkommene Zahlen sind.
\end{aufgabe}

\begin{loesung}
\begin{align*}
\sigma(6) &= 1 + 2 + 3 + 6 = 2 \cdot 6\\
\sigma(28) &= 1 + 2 + 4 + 7 + 14 + 28 = 2 \cdot 28
\end{align*}
\end{loesung}

\begin{aufgabe}[Aufgabe 5]
Bestimmen Sie drei Primzahlen der Form $2^k - 1$.
\end{aufgabe}

\begin{loesung}
$2^2 - 1 = 3$, \quad $2^3 - 1 = 7$, \quad $2^5 - 1 = 31$, \quad aber nicht $2^{11} - 1 = 2047 = 23 \cdot 89$
\end{loesung}

\begin{aufgabe}[Aufgabe 6]
Sei $2^k - 1$ eine Primzahl. Eine solche heißt Mersenne Primzahl.
\begin{enumerate}[label=\alph*)]
  \item Zeigen Sie, dass die Zahl $n = 2^{k-1}(2^k - 1)$ vollkommen ist.
  \item Bestimmen Sie eine vollkommene Zahl $> 28$.
  \item Zeigen Sie: Seien $m, n > 1$ ganze Zahlen, dann ist $2^{mn} - 1$ keine Primzahl.
\end{enumerate}
\end{aufgabe}

\begin{loesung}
\begin{enumerate}[label=\alph*)]
  % KORRIGIERT: Begründung ggT = 1 für die Multiplikativität fehlte
  \item Da $2^k - 1$ ungerade ist, gilt $\ggT(2^{k-1}, 2^k - 1) = 1$. Also
  $\sigma(n) = \sigma(2^{k-1}) \cdot \sigma(2^k - 1) = (2^k - 1) \cdot 2^k = 2\,n$
  \item $n = 2^4 \cdot (2^5 - 1) = 16 \cdot 31 = 496$ \quad ist vollkommen.
  \item \[ q^n - 1 = (q - 1)(1 + q + q^2 + \ldots + q^{n-1}) \]
  % KORRIGIERT: Begründung, dass beide Faktoren > 1 sind, fehlte
  Setzt man $q = 2^m$, folgt die Behauptung: Wegen $m > 1$ ist $q - 1 = 2^m - 1 > 1$, und wegen $n > 1$ ist $1 + q + \cdots + q^{n-1} > 1$.
\end{enumerate}
\end{loesung}

\subsection{Potenzen, Ordnung und Inverse in $\Zm{m}^*$}

\begin{aufgabe}[Aufgabe 7]
% KORRIGIERT: "3^1000 (mod 7)" als Wert, jetzt \bmod
Gesucht: $3^{1000}$ in $\Zm{7}$ bzw.\ $3^{1000} \bmod 7$
\end{aufgabe}

\begin{loesung}
% KORRIGIERT: Voraussetzung ggT(3, 7) = 1 für Fermat ergänzt
Nach Fermat ist $3^6 \equiv 1 \pmod 7$ (da $\ggT(3, 7) = 1$),\\
$1000 = 996 + 4 = 166 \cdot 6 + 4$\\
Also $3^{1000} = (3^6)^{166} \cdot 3^4 \equiv 3^4 = 81 = 77 + 4 \equiv 4 \pmod 7$
\end{loesung}

\begin{aufgabe}[Aufgabe 8]
Bestimmen Sie die Ordnung der Elemente $2, 3, 4, 5$ in $\Zm{17}^*$.\\
Welche dieser Elemente sind primitiv?\\
Bestimmen Sie die Inversen.
\end{aufgabe}

\begin{loesung}
% KORRIGIERT: Kontext "in Z_17" für die Gleichungen ergänzt
Alle Rechnungen in $\Zm{17}$:

\begin{tabular}{lcrcr}
$2,\ 4,\ 8,\ -1,\ -2,\ -4,\ -8,\ 1$ & $\Rightarrow$ & $\operatorname{ord}(2) = 8$ & $\Rightarrow$ & $2^{-1} = -8$\\
$3,\ -8,\ -7,\ -4,\ 5,\ -2,\ -6,\ -1,\ \ldots$ & $\Rightarrow$ & $\operatorname{ord}(3) = 16$ & $\Rightarrow$ & $3^{-1} = 6$\\
$4,\ -1,\ -4,\ 1$ & $\Rightarrow$ & $\operatorname{ord}(4) = 4$ & $\Rightarrow$ & $4^{-1} = -4$\\
$5,\ 8,\ 6,\ -4,\ -3,\ 2,\ -7,\ -1,\ \ldots$ & $\Rightarrow$ & $\operatorname{ord}(5) = 16$ & $\Rightarrow$ & $5^{-1} = 7$
\end{tabular}

\medskip
Primitiv sind 3, 5
\begin{align*}
2 \cdot 9 &= 18 \equiv 1 \pmod{17}\\
3 \cdot 6 &= 18 \equiv 1 \pmod{17}\\
4 \cdot 13 &= 52 \equiv 1 \pmod{17}\\
5 \cdot 7 &= 35 \equiv 1 \pmod{17}
\end{align*}
\end{loesung}

\textbf{Formel für $a^{-1}$ in $\Zm{m}^*$}\\
% KORRIGIERT: "a a^(phi(m)-1) = a^phi(m) = 1" ohne ggT-Voraussetzung und ohne mod m
Für alle $a$ mit $\ggT(a, m) = 1$ gilt (Satz von Euler) $a \cdot a^{\varphi(m)-1} = a^{\varphi(m)} \equiv 1 \pmod m$, also $a^{-1} \equiv a^{\varphi(m)-1} \pmod m$

\subsection{RSA}

\begin{aufgabe}[Aufgabe 9]
Alice will die Nachricht $a = 715$ (GO) verschlüsseln und als verschlüsselte Nachricht $c$ (Codewort) an Bob senden. Sie verabreden als Modul die Primzahl $m = 2803$ zu wählen und den öffentlichen Schlüssel $e = 113$ zu verwenden.
\begin{enumerate}[label=\alph*)]
  \item Mit welcher Formel verschlüsselt Alice das Wort $a$? Berechnen Sie das Codewort $c$. Hinweis: Zuerst $e$ im Binärsystem schreiben.
  \item Bob besitzt den privaten Schlüssel $d$. Wie (Formel) entschlüsselt Bob das Codewort $c$?
  \item Oscar erhält Kenntnis von $m$ und $e$. Welche Formeln kann er nutzen um den privaten Schlüssel $d$ zu berechnen? Geben Sie beide Lösungswege an und berechnen Sie $\varphi(\varphi(m))$.
\end{enumerate}
\end{aufgabe}

\begin{loesung}
1.\ $m = 2803$, $a = 715$
\begin{enumerate}[label=\alph*)]
  \item $e = 113$, binär $e = 1110001$
  \[
  % KORRIGIERT: "715^16 = 163" (richtig: 1653); "c = a^e" statt Kongruenz; "*" als Malzeichen
  c \equiv a^e = 715 \cdot 715^{16} \cdot 715^{32} \cdot 715^{64} \equiv 715 \cdot 1653 \cdot 2287 \cdot 2774 \equiv 708 \pmod{2803}
  \]
  % KORRIGIERT: "(mod m)" als Wert, jetzt \bmod
  \item $w = c^d \bmod m$
  \item 1.\ Lösungsweg
  \[
  % KORRIGIERT: "*" als Malzeichen (ganze Rechnung); Voraussetzung ggT(e, phi(m)) = 1 ergänzt
  e \cdot d \equiv 1 \pmod{\varphi(m)}. \qquad 113 \cdot d \equiv 1 \pmod{2802} \ (\varphi(2803) = 2802, \text{ da } 2803 \text{ prim}). \qquad \text{EA mit } 2802,\ 113
  \]
  Die Inverse von $e$ existiert, weil $\ggT(113, 2802) = 1$ (letzter Rest $\neq 0$ im EA ist $1$).
  \begin{align*}
  2802 &= 24 \cdot 113 + 90\\
  113 &= 1 \cdot 90 + 23\\
  90 &= 3 \cdot 23 + 21\\
  23 &= 1 \cdot 21 + 2\\
  21 &= 10 \cdot 2 + 1\\
  2 &= 2 \cdot 1
  \end{align*}
  \[
  Q = \begin{pmatrix} 24 & 1\\ 1 & 0 \end{pmatrix}
      \begin{pmatrix} 1 & 1\\ 1 & 0 \end{pmatrix}
      \begin{pmatrix} 3 & 1\\ 1 & 0 \end{pmatrix}
      \begin{pmatrix} 1 & 1\\ 1 & 0 \end{pmatrix}
      \begin{pmatrix} 10 & 1\\ 1 & 0 \end{pmatrix}
      \begin{pmatrix} 2 & 1\\ 1 & 0 \end{pmatrix}
    = \begin{pmatrix} 2802 & 1339\\ 113 & 54 \end{pmatrix}
  \]
  \[
  \det Q = 2802 \cdot 54 - 113 \cdot 1339 = 1
  \]
  \[
  \Rightarrow\quad d = -1339 \equiv 1463 \pmod{2802}
  \]

  2.\ Lösungsweg
  \begin{align*}
  & d \equiv e^{\varphi(\varphi(m)) - 1} \pmod{\varphi(m)} \quad (\text{Euler, da } \ggT(e, \varphi(m)) = 1).\\
  & \qquad \varphi(\varphi(2803)) = \varphi(2802), \quad 2802 = 2 \cdot 3 \cdot 467\\
  \Rightarrow\quad & \varphi(\varphi(2803)) = 2 \cdot 466 = 932\\
  \Rightarrow\quad & \varphi(\varphi(2803)) - 1 = 931 = 512 + 256 + 128 + 32 + 2 + 1
  \end{align*}
\end{enumerate}

Rechnung modulo 2802 (alle Gleichungen als $\equiv \pmod{2802}$ zu lesen)
\begin{align*}
e &= 113\\
e^2 &= 1561\\
e^4 &= 1783\\
e^8 &= 1621\\
e^{16} &= 2167\\
e^{32} &= 2539\\
e^{64} &= 1921\\
e^{128} &= 7\\
e^{256} &= 49\\
e^{512} &= 2401
\end{align*}
\[
d = 113 \cdot 1561 \cdot 2539 \cdot 7 \cdot 49 \cdot 2401 \equiv 1463 \pmod{2802}
\]

\medskip
Probe:
\[
% KORRIGIERT: "a = c^d = ..." statt Kongruenz; "822 161" ohne Malzeichen
a \equiv c^d = 708^{1463} \equiv 708 \cdot 2330 \cdot 2292 \cdot 1957 \cdot 951 \cdot 822 \cdot 161 \cdot 2323 \equiv 715 \pmod{2803}
\]

\[
1463 = 1024 + 256 + 128 + 32 + 16 + 4 + 2 + 1
\]
\begin{align*}
c^1 &= 708 \qquad \text{alles modulo } 2803\\
c^2 &= 2330\\
c^4 &= 2292\\
c^8 &= 442\\
c^{16} &= 1957\\
c^{32} &= 951\\
c^{64} &= 1835\\
c^{128} &= 822\\
c^{256} &= 161\\
c^{512} &= 694\\
c^{1024} &= 2323
\end{align*}
\end{loesung}
